%============================================================
\chapter{実証の設計と調査計画：観測への接続}
\label{ch:design}
%============================================================

\section{何を実証するのか}

第\ref{part:theory}部は量と命題を定義し、第\ref{part:catalog}部は制約下の部分空間を演繹した。
いずれも観測を用いていない。本部以降は、これらを観測に接続する。

接続の対象は二つに分かれる。

\begin{center}
\begin{tabularx}{\textwidth}{@{}lXl@{}}
\toprule
対象 & 内容 & 問い \\
\midrule
第\ref{part:theory}部の量 & $\kappa$、$\CCC$、$g^\star$、$\phi$、$M(t)$、$\mathcal{G}$ など & 実際にどの値をとるか \\
第\ref{part:catalog}部の演繹 & 類型のふるい分け、代替関係、無人化、担保制約 & 前提の下で成立するか \\
\bottomrule
\end{tabularx}
\captionof{table}{実証が接続する二つの対象。}
\label{tab:empirical-targets}
\end{center}

前者は測定であり、後者は検証である。
両者は必要な準備が異なるため、以下では段階に分けて扱う。

\section{通過すべき五段階}
\label{sec:five-stages}

一つの量、または一つの命題を観測に接続するまでに通過すべき段階を五つに分ける。
この整理は本稿のものである。

\begin{center}
\begin{tabularx}{\textwidth}{@{}llXl@{}}
\toprule
 & 段階 & 問い & 通過しない場合に要するもの \\
\midrule
1 & 操作化 & その量に単位と観測対象が定まるか & 定義の精密化 \\
2 & 母集団 & 何を数え上げる集合とするか & 対象の選び直し \\
3 & 測定 & 必要な粒度で値が取れるか & 取得経路の開拓 \\
\midrule
4 & 照合 & 既存の集計値と整合するか & --- \\
5 & 検証 & 事前に固定した含意が支持されるか & --- \\
\bottomrule
\end{tabularx}
\captionof{table}{理論を観測に接続するまでの五段階。}
\label{tab:five-stages}
\end{center}

\textbf{前三者は直列の前提条件である。} いずれかを欠けば後の段階に進めない。
\textbf{後二者は取得した値の使い方であり、段階の続きではない。}
両者は事前登録の有無で分かれる。

段階を分ける利点は、詰まった箇所によって要する作業が異なる点にある。
段階1の穴はデータをいくら集めても埋まらず、段階3の穴は定義をいくら精密にしても埋まらない。
\textbf{どこで詰まっているかを先に判定しないと、労力の配分を誤る。}

\section{段階ごとの確認事項}

\subsection{段階1：操作化}

確認するのは三点である。

\begin{enumerate}[label=(\arabic*)]
\item その量に\textbf{単位}が定まるか
\item \textbf{観測対象}が何であるかが定まるか。企業か、$\Phi$ か、識別子か
\item 命題であれば、\textbf{符号が一意に定まるか}
\end{enumerate}

第\ref{part:theory}部の定義がそのまま操作化になっている量と、そうでない量がある。
前者については第\ref{ch:next}節に $\Phi$ ごとに記録すべき変数を列挙する。

後者、すなわち定義は与えたが単位や観測対象が定まらない量については、
測定に進む前にその旨を確定させる。
\textbf{この段階の不備は、データを追加しても検出できない。}
値が得られてしまうため、何を測ったのかが分からないまま結果が出る。

(3) は命題に固有である。理論が符号を一意に定めない場合、
どちらの結果が出ても理論と整合してしまい、検定が成立しない。

\subsection{段階2：母集団}

確認するのは三点である。

\begin{enumerate}[label=(\arabic*)]
\item 基準時点、範囲、下限を定めたか
\item 標本に入らない対象は何か。それは調べたい軸と相関するか
\item 停止・退出した対象が標本に残るか
\end{enumerate}

(2)(3) が本部の中心的な主題であり、第\ref{ch:method}章がこれを扱う。
本稿の枠組みでは (3) が特に重い。
第\ref{sec:objective}節の倒産時刻 $\tau$ は $\Phi$ に依存するため、
$\Phi$ を推定しようとすると標本の構成そのものが推定対象の関数になる。

\subsection{段階3：測定}

確認するのは三点である。

\begin{enumerate}[label=(\arabic*)]
\item その量がどこかで測定されているか
\item 測定されている場合、必要な粒度で取得できるか
\item 取得できない場合、障害はどの種別か
\end{enumerate}

(3) の判定が調査設計を左右する。
障害の種別によって、労力の投入が前進を生む場合と何も生まない場合が分かれるためである。
種別の分類は調査を経て構成されるため、第\ref{part:obstacles}章に置く。

\begin{remark}[本章が分類そのものを与えない理由]
\label{rem:taxonomy-placement}
段階3の判定には障害の分類が要るが、その分類は実証を経て構成されている。
本章に前置きすれば、まだ得られていない知見を先に置くことになる。

\textbf{本章が与えるのは確認すべき段階であり、通過できなかった原因の分類ではない。}
前者は設計の手続きであり、後者は調査の結果である。
両者は独立に述べられる。段階は「どこで詰まったか」を、
分類は「なぜ詰まったか」を与える。
\end{remark}

\subsection{段階4と5：取得した値の使い方}

段階3までを通過すれば値が得られる。その使い方が二つに分かれる。
次節で扱う。

\section{照合と検証の区別}
\label{sec:comparison-vs-test}

演繹の帰結と既存の集計値を事後に見比べることを\emph{照合}と呼ぶ。
事前に固定した含意をデータで判定することを\emph{検証}と呼ぶ。

\begin{center}
\begin{tabularx}{\textwidth}{@{}lXX@{}}
\toprule
 & 照合 & 検証 \\
\midrule
含意の固定 & 事後 & 事前 \\
母集団の宣言 & 不要 & 必要 \\
一致しない場合 & 別の要因により説明できる & 含意が棄却される \\
反証可能性 & \textbf{持たない} & 持つ \\
\bottomrule
\end{tabularx}
\captionof{table}{照合と検証の区別。}
\label{tab:comparison-vs-test}
\end{center}

\textbf{照合は反証されない。} 一致しなければ別の要因が働いたと述べることができ、
その説明を排除する手段が標本の中にない。したがって照合は理論を支持も反証もしない。

それでも照合には用途がある。
演繹が現実と大きく食い違っていないことの確認であり、
次に何を検証するかの選定に使える。
\textbf{確認としての価値はあるが、証拠としての価値はない。}

危険は別のところにある。

\begin{remark}[照合は段階を飛ばせる]
\label{rem:comparison-skips-stages}
照合は既存の集計値を用いるため、\textbf{自分の母集団を宣言しなくても成立する}。
すなわち段階2と3を経ずに段階4に到達できる。

これが照合と検証が混同される経路である。
段階1から3を通過した上で含意が支持された場合と、
段階2と3を飛ばして既存の集計値と方向が一致した場合とでは、
得られている根拠の強さが桁で異なる。
外形上はどちらも「理論と整合した」と書ける。
\end{remark}

第\ref{ch:solo}章で族3の復活を年齢別の起業分野と照合したのは、この形にあたる。
同節は事後的な照合であり因果を示すものではない旨を明記している。

以上より、次の規律を置く。

\begin{quote}
\textbf{照合を行うときは、それが照合であることを明示する。
事前に固定した含意の代わりに用いない。}
\end{quote}

これは第\ref{ch:method}章の「有名企業は例示に用い、推論には用いない」と同型の規則である。
いずれも、\textbf{得られた材料が持つ根拠の強さを超えて用いないこと}を要求する。

\section{本稿における段階の対応}

各段階を本稿のどこで扱うかを示す。

\begin{center}
\begin{tabularx}{\textwidth}{@{}llX@{}}
\toprule
段階 & 扱う箇所 & 内容 \\
\midrule
1 操作化 & 第\ref{part:theory}部、第\ref{ch:next}節 & 定義と、$\Phi$ ごとに記録すべき変数 \\
2 母集団 & 第\ref{ch:method}章 & 選択バイアスの分解と対処方針 \\
3 測定 & 第\ref{ch:instantiation}章、第\ref{part:obstacles}章 & 測れた量と測れなかった量、障害の分類 \\
4 照合 & 本章 & 第\ref{sec:comparison-vs-test}節の規律 \\
5 検証 \& 第\ref{ch:next}節、第\ref{part:results}部 \& 事前登録の手続きと、22の含意の判定 \\
\bottomrule
\end{tabularx}
\captionof{table}{五段階と、本稿においてそれを扱う箇所。}
\label{tab:stage-mapping}
\end{center}

第\ref{ch:instantiation}章を第\ref{part:results}部の冒頭に置いたのは、
段階3の結果を一覧するためである。
測定できなかった量を明示することが同章の主目的であり、
それは段階1から3のどこで詰まったかの記録にあたる。

\begin{remark}[段階は量ごとに異なる]
\label{rem:stage-per-quantity}
到達段階は理論の全体について一つに定まるのではなく、量ごと・命題ごとに異なる。
同一の章の中で、段階5まで到達した量と段階1で止まっている量が併存する。

したがって「本稿の実証はどこまで進んだか」という問いには単一の答えがない。
到達段階は量ごとに述べるほかない。
\end{remark}

%============================================================

%---- 調査計画 ----
\label{ch:next}

%============================================================

\section{調査の優先順位}

対象を族2 と族5 に絞る。選定は演繹で導ける。

\paragraph{$\phi_{\mathrm{prod}}$ は $\Phi$ に依存しない。}

式\eqref{eq:surplus}の三分解のうち $\phi_{\mathrm{prod}}=(p^{\ast}-c(\delta))\delta$ は
限界費用と競争価格で決まり、履行から決済への写像の形を含まない。
同一の $f$ と同一の限界費用に対して $\Phi$ だけを差し替えても、
変わるのは $\kappa$ の符号と必要な運転資本であって生産余剰ではない
（表\ref{tab:same-f-different-phi}）。

\textbf{$\Phi$ の形が余剰に効くのは残る二項においてである。}
枠組みを検証するには、この二項をそれぞれ単独に含む $\Phi$ を選べばよい。

\paragraph{二項は別の自由度から生じる。}

$\phi_{\mathrm{cog}}=p\cdot\E[\dbar-\delta]$ が正となりうるのは、
契約が上限 $\dbar$ を明示する場合、すなわち $\pi$ が $\delta$ を参照しない場合に限られる
（注\ref{prop:breakage}）。これは自由度(5)である。
加えて、乖離を取引頻度 $\nu$ に結びつけるには自由度(6)の反復性を要する。
単発では $\dbar-\delta$ が一回の実現でしかなく分布を持たないため、
$\nu$ との関係を問えない。
\textbf{自由度(5)と(6)がともに立つ族は族2}であり、
第\ref{part:catalog}部はこれを認知余剰の主要な所在として定めている。

$\phi_{\mathrm{barg}}=(p-p^{\ast})\delta$ はゼロサムの移転であり、
その配分は $\pi=h(y)$ の $y$ が検証可能かどうかで決まる。これは自由度(4)である。
表\ref{tab:family-assignment-order}の判定順序において、
\textbf{自由度(4)で判定される族は族5のみ}である。

\paragraph{二族は象限図の対角にある。}

自由度(1)と(2)が支配的であり、この二軸が図\ref{fig:quadrant}の四分割を与える。
族2 は前受・定額型（$\pi$ が $\delta$ と独立、乖離しうる）、
族5 は後払・連動型（$\pi$ が $\delta$ に連動、乖離しない）であって、
\textbf{二軸の値がいずれも逆になる}。
どちらか一族だけを選べば、動かせるのは二軸のうち片方だけである。

\paragraph{何を検証できるか。}

族2 は認知余剰（式\eqref{eq:cog}）と容量制約の命題を直接検証できる。
特に、同一の定額会員制でありながら容量制約の有無で二分される点は、
外形的に似た $\Phi$ が異なる耐久性を持つことの最も明快な事例になる。

族5 は可測性条件（第\ref{sec:info}章）を直接検証できる。
成果報酬とコストプラスが逆向きのリスク移転であるという主張は、
契約形態と観測可能性の対応として検証可能である。

\paragraph{並べる理由。}

二族は制度的な根拠が異なるため、データの得られ方も異なる。
族2 は資金決済法の登録制により全数フレームが存在し、
族5 は職業安定法の情報提供義務により個社の届出が公開されている。
\textbf{前者は網羅的な枠を与え、後者は個体の変数を与える。}

一族だけを見ると、検証が進まなかったときに、
その原因を\textbf{フレームの不備に帰すべきか変数の不備に帰すべきか決まらない}。
根拠の異なる二族を並べれば、どちらの側で詰まったのかを切り分けられる。
二族を対象とする理由はここにもある。

\section{記録すべき変数}

各 $\Phi$ について、最低限次を記録する。

\begin{center}
\begin{tabularx}{\textwidth}{@{}lX@{}}
\toprule
変数 & 内容 \\
\midrule
$\operatorname{sgn}\kappa$ & 決済が履行に先行するか後行するか \\
$\CCC$ & 定常状態が成立する場合のみ。成立しない場合はその旨を記録 \\
依存形式 & $c$ / $p\delta$ / $F+p\delta$ / $h(y)$ のいずれか \\
$\dbar$ と $\delta$ の分離 & 分離の有無、および乖離の推定値 \\
容量制約 & 有無、および制約が拘束的かどうか \\
$\nu$ & 取引頻度（式\eqref{eq:decay}） \\
可測性 & 努力・結果のいずれが検証可能か \\
支払主体 & 受益者本人か、第三者か \\
$\phi$ の推定内訳 & $\phi_{\mathrm{prod}}$ / $\phi_{\mathrm{barg}}$ / $\phi_{\mathrm{cog}}$ の定性的配分 \\
\bottomrule
\end{tabularx}
\captionof{table}{$\Phi$ ごとに最低限記録すべき変数の一覧。}
\label{tab:variables-to-record}
\end{center}

\section{事前登録すべき事項}

第\ref{part:method}部の結論に従い、調査開始前に以下を確定し記録する。

\begin{enumerate}[label=(\arabic*)]
\item 母集団の定義（基準時点、地理的範囲、業種範囲、規模下限）
\item 定常状態とみなす閾値（$r$ の安定性の判定基準）
\item 失敗事例の収集方法と、収集不能な範囲
\item APC のうちゼロと仮定する効果、およびその根拠
\item 一次資料と回顧的記述の区別の基準
\end{enumerate}

これらを事後に決めると、第\ref{sec:survivorship}節および第\ref{sec:disclosure}節の
選択が結論に混入する。

